%% ---------------------------------------------------------------------------
%% 1. Introduction
%% ---------------------------------------------------------------------------
\section{Introduction}

As a node-based numerical method, the meshfree methods enable an efficient and straightforward way to construct arbitrary order smooth shape function \cite{ref1,ref2,ref3,ref4,ref5,ref6,ref7} and execute post process without full-connected elements like in finite element methods \cite{ref8}.
In weak form based meshfree methods, Galerkin meshfree methods invoke Galerkin weak form of equilibrium equations thus lower the continuity requirements of approximation, which greatly improves the stability and accuracy \cite{ref1,ref5}.
Consequently, Galerkin meshfree methods have attracted lots of attention among scholars and developed into various kinds of approaches, such as Element-free Galerkin method \cite{ref9,ref10}, reproducing kernel particle method \cite{ref11}, max-entropy meshfree method \cite{ref12}, smoothed particle Galerkin method \cite{ref13}, among others \cite{ref14,ref20}.

Additionally, due to the arbitrary order smooth approximants, it is a desirable way to apply the Galerkin meshfree methods on the higher order problems like Kirchhoff--Love plate theory, which requires at least $C^1$ continuity of shape function for Galerkin weak form.
Specifically, Krysl et al. \cite{ref21} adopted element-free Galerkin method to analyze the thin plates.
Subsequently, Liu et al. \cite{ref22,ref23} theoretically investigate the performance of element-free Galerkin method in static and dynamic analysis of thin plates.
Petrov--Galerkin meshfree method is applied in static and dynamic analysis by Gu and Liu \cite{ref24}.
By combining triangular integration cell with RK approximation, Simkins et al. \cite{ref25} proposed reproducing kernel element method and assessment its performance in solving Kirchhoff--Love plate problems.
With the expectation of improving efficiency and conveniently imposing rotation boundary condition, Wang and Chen \cite{ref26} introduced rotation variable into reproducing kernel approximation, and presented a Hermite reproducing kernel approximation for thin plates analysis.
By introducing reproducing polynomial particle approximation, Oh et al. \cite{ref27} proposed a meshfree particle method for thin plates, which has high order reproducing property and the Kronecker delta property.
Based on linear basis function, Deng et al. \cite{ref28} proposed a gradient smoothing Galerkin meshfree method to eliminate the necessity of $C^1$ continuous approximation.

It should be noted that meshfree shape function usually possess rational nature, which brings severe difficulty towards accurate numerical integration \cite{ref29,ref30}.
The integration constraints which has been numerically proved play a crucial role in convergence performance of Galerkin meshfree methods \cite{ref31,ref32,ref33,ref34,ref35}, was first proposed and validated by Chen et al \cite{ref31,ref32}.
In order to satisfy the integration constraints or realize integration consistency, lots of approaches have been presented and applied in many problems \cite{ref33,ref34,ref35,ref36,ref37,ref38,ref39,ref40,ref41,ref42,ref43,ref44}.
More specifically, Chen et al. \cite{ref31,ref32} replaced the derivate of shape function by strain smoothing technology such that the linear integration constraints have been satisfied.
By solving sets of equations generated by quadratic integration constraints, Duan et al. \cite{ref35,ref40} proposed quadratically consistent meshfree integration method, which are further developed and applied in strain gradient elasticity.
Wang et al. \cite{ref42} extended SCNI into quadratic meshfree formulation and proposed sub-domain stabilized conforming integration scheme in order to meet the quadratic integration constraints.
Based upon reproducing kernel approximation, Wang and Wu \cite{ref43} presented an inherently consistent reproducing kernel gradient smoothing integration algorithm (RKGSI), which automatically meet arbitrary order of basis function.
Thereafter, Du et al. \cite{ref45} presented a unified reproducing kernel gradient smoothing approach as extension and application of RKGSI in higher order problems.

Regarding to error estimates of Galerkin meshfree methods, Liu et al. \cite{ref46} measured the convergence rate of MLS reproducing kernel methods (MLSRKPM) by a dilation parameter of window function.
Salehi and Dehghan \cite{ref47} conducted the error estimated for a class of time-dependent partial differential equations studied by the proposed MLS radial reproducing polynomial meshless method (MLSRRPM).
Han and Meng \cite{ref48} introduced a theoretical analysis of the reproducing kernel particle method (RKPM) by interpolation error estimates.
Osorio and French \cite{ref49} theoretically validated the advantages of proposed stabilized diffuse Galerkin method (SDGM) by providing error estimates.
Among them, the importance of integration consistency has not been theoretically proved although lots of numerical experiments have well shown the desirable convergence performance of methods with integration consistency.

Recently, Wu and Wang \cite{ref50} took the effect of numerical integration into consideration, providing a valuable error analysis of Galerkin meshfree methods for second order problems like potential problems and elasticity problems.
As an extension of the previous work, we focus on the fourth order Kirchhoff--Love plate problems, investigating the crucial role of integration constraints and meshfree approximations.
Firstly, we estimate the convergence rate of $H_2$ error of Galerkin meshfree methods for thin plates, which is strongly related to the order of integration consistency and meshfree approximation.
Then, with the help of the closed graph theorem \cite{ref51} and Lagrange interpolants, we further investigate the $L_2$ and $H_1$ error estimates by dividing them into two parts, which is related to the standard error of meshfree approximation and error of orthogonality, respectively.
Subsequently, in light of constrains on the order selection of Lagrange interpolation from meshfree approximation, the error estimate for first part is further divided into two cases, the order of meshfree approximation is greater than or equal to three ($p\ge3$), and equal to two ($p=2$).
It turns out that, only when the order of meshfree approximation satisfy the minimum requirement ($p\ge3$) and integration schemes possess corresponding order of integration constraints, can the Galerkin meshfree methods achieve optimal convergence performance.
The numerical experiments well demonstrate the presented theoretical results.

The remainder of this paper is organized as follows.
Section \ref{sec:formulation} introduces the discrete Galerkin meshfree formulation for Kirchhoff--Love plate, and the reproducing kernel gradient smoothing integration scheme is described.
Section \ref{sec:error} is devoted to the error estimates of the Galerkin meshfree formulation in $H_2$, $H_1$ and $L_2$ norms.
Numerical examples are examined in Section \ref{sec:numerical} and the conclusions are drawn in Section \ref{sec:conclusion}.
